The paper belongs to the literature on the real effects of financial markets \citep{BondEdmansGoldstein2012}.
The closest antecedent is \citet{DowGoldsteinGuembel2017}, where a firm's investment decision feeds back into
the incentive to produce information about its stock. Here the sale rule divides acquisition surplus between a
traded claim and a buyer that decides whether to enter. \citet{EdmansGoldsteinJiang2015} show that corrective
real decisions can discourage trade on bad news under rational pricing. Here the real decision is the buyer's
entry itself.

Firms learn from peers' stock prices \citep{FoucaultFresard2014}, and target prices affect takeover activity
\citep{EdmansGoldsteinJiang2012}. In this model, a prospective challenger learns about its own acquisition value
from the target's price before it commits.\footnote{Bidders also learn from announcement returns when they decide
whether to complete a deal \citep{Luo2005}. That decision comes after the entry margin of this paper.}

The auction-entry literature makes the bidder pool endogenous
\citep{Samuelson1985,McAfeeMcMillan1987,LevinSmith1994,GentryStroup2019}. Under selective entry, a bidder's
pre-entry information indexes the distribution of its value, and the payment of the entry cost reveals the value
\citep{RobertsSweeting2013}.\footnote{One limit of this framework is a bidder that knows only the common value
distribution before entry \citep{LevinSmith1994}; the other is a bidder that already knows its value
\citep{Samuelson1985}. Selective entry also affects the choice of sale procedure \citep{RobertsSweeting2013}, and
auction formats affect bidders' incentives to acquire information \citep{Persico2000}.} Here the pre-entry
information is the target's stock price. That price is public, and it moves with the orders of an investor who
observes the challenger's quality class. I keep the direct deterrence force of the auction-entry literature and add a
market that operates before preparation. The informed trader is outside the auction and trades a payoff that differs
from the entrant's profit.

The investor in this paper holds no stake, seeks no control, has no role in tendering, and earns only the trading
residual. An acquirer can instead use liquidity trade to conceal share purchases before a bid \citep{KyleVila1991},
large holders trade on their information \citep{Edmans2009}, and arbitrageurs' positions can affect tendering
\citep{CornelliLi2002}.\footnote{Large holders also monitor or intervene as prices feed back to firm decisions
\citetext{\citealp[Section~II.E]{EdmansGoldsteinJiang2015}; \citealp{Maug1998}; \citealp{BackEtAl2018}}.
\citet{GoldsteinGuembel2008} study manipulation through the real decision that responds to the price. The model
abstracts from toeholds and dispersed-shareholder free riding \citep{BulowHuangKlemperer1999,GrossmanHart1980}.
Recent work studies how bidders learn about their own values and about competitors \citep{PernoudGleyze2026},
post-auction feedback in security-payment design \citep{LiuBernhardt2022}, and bidder-pool choice with correlated
values \citep{CarlinEtAl2026}. In takeovers, \citet{BettonEtAl2014} study negotiations with stock-market feedback,
and \citet{LinMaYangZhu2025} model payment choice and trade within an initiated deal. I hold cash consideration fixed
and study entry.}

Section~\ref{sec:model} sets out the model, and Section~\ref{sec:payoffs} derives the returns to information.
Section~\ref{sec:results} states the main results. Sections~\ref{sec:research} to~\ref{sec:empirical} cover research,
robustness, sale terms and empirical implications, and Section~\ref{sec:conclusion} concludes.
Appendix~\ref{app:proofs} states the results of Sections~\ref{sec:robustness} to~\ref{sec:empirical} in full and
proves every result.
